\scanpage{4}
et comme $\dim A_0\leq\dim A'_0$ par ce qui précède, on a égalité partout ; \unsure{ceci|ce} prouve en particulier que \uline{$A'_0$ est simple} ], et aussi
\[
\dim H^{1}(A_0,\underline O_{A_0})=n,\quad\text{d'où la structure totale de }H^{\unsure{x|*}}(A_0,\underline O_{A_0})\unsure{\big]|.}
\]

Comme on vient de voir que $A'_0$ est simple, on voit que pour tt $x'\in A'_0$, $\exists\,x\in A_0$ au dessus et on aura nécessairement $\underline O_{\deleted{\unsure{X}}A'_0,x'}\unsure{\text{---}|\longrightarrow}\underline O_{A_0,x}$ \uline{injectif}.

Or on a le

\uline{Lemme} $X$ et $Y$ de type fini sur $S$ loc. noeth. \inserted{\unsure{avec|soient} $F,G$ \unsure{quacoh.|quacoh|quasi-coh.} sur $X,Y$} $\varphi:X\to Y$ un $S$-morphisme, $s\in S$, \inserted{$G\to F$ un $\varphi$-morphisme}, \deleted{\unsure{on suppose}} $y\in Y$ au dessus de $y$, on suppose a) $\deleted{\unsure{G}} \unsure{F|G}$ est plat sur $\unsure{S_y|S,|S}$ b/ $G$ coh.

% slanted
\margin{mettre aussi le cas où l'hypo. noeth. \deleted{\unsure{de finitude|de Noether}} \unsure{peut-être|peut être} sera remplacée par l'hypothèse que $Y$ \deleted{\unsure{réduit}} est local à radical nilpotent \ldots}

b) Soit $\varphi_s:X_s\to Y_s$ induit par $\varphi$, \deleted{\unsure{alors}} tout \unsure{élémt|élém$^{\mathrm t}$|élément} de $G_s$ dont l'image inverse en les $x\in\varphi_s^{-1}(y)$ s'annule est \inserted{nulle} \deleted{$G_s\to\varphi_{s,*}(F_s)$ est injectif en $y$.} \uline{Alors}

$G$ est plat sur $S$ en $y$.

On utilise le critère par les gradués associés.

On applique ceci au $\varphi:A\to A_{\deleted{\unsure{0}}}'$ et un pt quelconque de $A'$, $F=\underline O_A$, $G=\underline O_{A'}$.

Cela prouve que $A'$ \unsure{est|$\to$} plat sur $S$ \uline{cqfd}.
% archive sheet number: 3

